\chapter{Noninvasive Ventilation}

\section*{Overview}
Noninvasive ventilation (NIV) is the delivery of positive-pressure ventilatory support through a mask without an endotracheal tube. It is used to support patients with acute or chronic respiratory failure while preserving airway defenses and speech.

\section*{Alternative Names}
NIV, bi-level positive airway pressure (BiPAP), continuous positive airway pressure (CPAP), mask ventilation, noninvasive positive pressure ventilation

\section*{Principle}
A face or nasal mask delivers pressurized gas to the upper airway, providing inspiratory support and/or continuous positive end-expiratory pressure. This unloads the respiratory muscles, recruits collapsed alveoli, and splints the airway open, improving ventilation and oxygenation while avoiding the complications of intubation.

\section*{History}
Continuous positive airway pressure for obstructive sleep apnea was introduced in 1981, and noninvasive ventilation for acute respiratory failure became established in the 1980s and 1990s. It is now a standard first-line therapy for acute hypercapnic respiratory failure and cardiogenic pulmonary edema.

\section*{Why This Test or Procedure Is Ordered}
Ordered for acute exacerbations of COPD with respiratory acidosis, cardiogenic pulmonary edema, and acute asthma or pneumonia when selected, as well as for hypoxemic respiratory failure and immunocompromised patients to avoid intubation. It is also used chronically for sleep apnea, obesity hypoventilation syndrome, and neuromuscular disease.

\section*{Is This a Primary or Secondary Test or Procedure}
A primary treatment for selected acute and chronic respiratory failure, and an adjunct that may avoid or defer the need for invasive mechanical ventilation.

\section*{How This Test or Procedure Is Performed}
A well-fitted mask is secured over the nose or nose and mouth, and a ventilator or portable device delivers pressure with settings adjusted for inspiratory and expiratory positive airway pressure, fraction of inspired oxygen, and rate. The patient is coached on mask tolerance and synchronization, and response is monitored with pulse oximetry, blood gases, and clinical assessment.

\section*{What Preparation Is Needed}
Patient education, careful mask fitting and securing of the airway, and baseline blood gas and oximetry measurements. The patient should be awake, cooperative, and able to protect the airway.

\section*{Contraindications and Precautions}
Contraindications include respiratory or cardiac arrest, inability to protect the airway, copious secretions, vomiting, severe facial trauma or surgery, and hemodynamic instability. Precautions include gastric insufflation with aspiration risk, pressure-related skin injury, and the need for early conversion to invasive ventilation if NIV fails.

\section*{Risks}
Facial pressure sores, skin breakdown, eye irritation, gastric distention and aspiration, air leak, and failure of therapy with respiratory deterioration.

\section*{What Happens After the Test or Procedure}
The patient is monitored closely, and therapy is titrated based on blood gases and respiratory effort, with early reassessment within one to two hours. NIV is weaned gradually as the patient improves and discontinued when stable on supplemental oxygen.

\section*{Understanding Your Results}
Improvement in pH, PaCO2, oxygenation, and respiratory rate indicates a successful response, whereas worsening acidosis, inability to tolerate the mask, or hemodynamic decline signals NIV failure requiring intubation.

\section*{Normal Results}
Restoration of near-normal pH, ventilation, and oxygenation with comfortable breathing and a reduced work of breathing.

\section*{Follow-up Tests and Procedures}
Repeat blood gases and oximetry during titration, and a definitive plan for home NIV if chronic indications are present. Patients who fail NIV proceed to invasive mechanical ventilation and further evaluation of the cause of respiratory failure.

\section*{References}
\begin{enumerate}
  \item MedlinePlus. Noninvasive ventilation. U.S. National Library of Medicine; 2025.
  \item Current Medical Diagnosis and Treatment. McGraw-Hill; 2025.
\end{enumerate}

\clearpage
